\section{Multiplicity perturbation and proof of the main theorem}
\label{sec:multiplicity-endgame}

\begin{lemma}[Generic affine perturbation at fixed height]
\label{lem:generic-affine}
Fix $T$, the good horizontal contours, and the admissible $T^{-100}$ intervals
for the packet centres. Given any real neighbourhood $U_0$ of $(0,0)$, there
are parameters $(\eta,\varepsilon)\in U_0$, chosen compatibly with the Hardy
conjugation symmetry, such that
\[
H(z)=F(z)+\eta z+\varepsilon
\]
has only simple zeros and simple stationary zeros in $\Omega_T$, and $H,H'$ have
no common zero. Choose one packet centre per stationary node so that
Lemma~\ref{lem:evaluation-rank} gives an evaluation isomorphism. On this packet
space define $V_{\rm core}$ by Lemma~\ref{lem:stable-core-frame},
$V_{\rm mean}$ by \eqref{eq:mean-map}, $V_{\rm end}$ by
\eqref{eq:end-space}, and $\Pret$ by \eqref{eq:retained-space}. Then
\eqref{eq:retained-codim} holds. Write
\[
\Gfr_T^{\rm ret}:=\Gfr_T|_{\Pret}.
\]
Hypothesis~\ref{hyp:edge-regular-bounds} is therefore applied to the
original function $F$ on exactly the same finite packet space on which the
zero-side evaluation of $H$ is performed. Moreover the perturbation can be
chosen so small that, on this selected $\Pret$,
\begin{equation}
\boxed{
\left\|
(\Gfr_T^{\rm ret})^{-1/2}(A_H-A_F)
(\Gfr_T^{\rm ret})^{-1/2}
\right\|_{\op}<\frac{0.98\lambda_T}{8}.
}
\label{eq:perturb-small}
\end{equation}
\end{lemma}

\begin{proof}
Fix $T$. The chosen boundary contains no zero of $F$ or $F'$. Intersecting
$U_0$ with a sufficiently small convex real neighbourhood of the origin gives
$U\subset U_0$ on which $H$ and $H'$ remain nonvanishing there for every point
of the line segment from $(0,0)$ to $(\eta,\varepsilon)\in U$. Hence homotopy
invariance of the argument principle gives
\[
N_{\rm good}(H')=N_{\rm good}(F')=:m
\qquad ((\eta,\varepsilon)\in U).
\]
There are only finitely many admissible $m$-coordinate subsets, and each centre
ranges over a compact $T^{-100}$ interval. The corresponding reference Grams
are continuous and positive definite by Section~\ref{sec:ledger}; hence
\[
\gamma_T:=\inf_{\rm admissible\ centres}\lambda_{\min}(\Gfr_T)>0.
\]
Restriction to $\Pret$ can only increase the least eigenvalue. Uniform
continuity of the contour integrand on the same compact centre set allows $U$
to be shrunk so that
\begin{equation}
\sup_{(\eta,\varepsilon)\in U}
\sup_{\rm admissible\ centres}\|A_H-A_F\|_{\op}
<\frac{0.98\lambda_T\gamma_T}{8}.
\label{eq:uniform-fixedT-perturbation}
\end{equation}
No asymptotic lower bound for $\gamma_T$ is required.

It remains to impose genericity inside this fixed neighbourhood. Since
$H'=F'+\eta$ and $H''=F''$, a multiple stationary zero would satisfy
\[
F''(z)=0,\qquad \eta=-F'(z).
\]
The zeros of $F''$ in the fixed rectangle are finite, so only finitely many real
$\eta$ are forbidden. Choose a nonzero real $\eta$ in $U$ outside this set.
The stationary zeros $z_1,\dots,z_m$ are then simple. For this fixed $\eta$, a
common zero of $H,H'$ at $z_j$ would force
\[
\varepsilon=-F(z_j)-\eta z_j.
\]
Only finitely many real $\varepsilon$ are therefore forbidden. Choose a real
$\varepsilon$ in $U$ outside them. Then $H$ has only simple zeros, $H'$ has only
simple zeros, and $H,H'$ have no common zero. Real parameters preserve
$H(\bar z)=\overline{H(z)}$.

Lemma~\ref{lem:packet-dimension} now supplies at least $m$ admissible packet
intervals after the endpoint collars are removed. Choose $m$ distinct intervals
and then, by Lemma~\ref{lem:evaluation-rank}, centres for which the evaluation
map is an isomorphism. Lemma~\ref{lem:coordinate-restriction} preserves all
preceding source, frame, rank, and relative-remainder estimates on this
subfamily. Finally \eqref{eq:uniform-fixedT-perturbation} and
$\lambda_{\min}(\Gfr_T^{\rm ret})\ge\gamma_T$ give
\eqref{eq:perturb-small}. The right-edge source evaluation remains that of the
original function $F$; $H$ is used only for the zero-coordinate representation
and this finite-matrix comparison.
\end{proof}

\begin{lemma}[Off-axis zero count for the generic perturbation]
\label{lem:generic-zero-count}
Let $D(H)$ be the number of nonreal conjugate zero pairs of the straightened
generic $H$ in the chosen good rectangle $\Omega_T$. Then
\begin{equation}
\boxed{D(H)=B(H)+P(H')+o(N_T).}
\label{eq:D-H}
\end{equation}
\end{lemma}

\begin{proof}
Choose the perturbation in Lemma~\ref{lem:generic-affine} within the boundary
neighbourhood fixed in its proof. Homotopy invariance of the argument principle
gives exact equality of the zero counts of $F,H$ and of $F',H'$ inside the good
rectangle:
\[
N_{\rm good}(H)=N_{\rm good}(F),
\qquad
N_{\rm good}(H')=N_{\rm good}(F').
\]
By the good-rectangle form \eqref{eq:z1-zeta-good-decrement}, equivalently its
$(F',F)$ version, it follows that
\begin{equation}
N_{\rm good}(H')-N_{\rm good}(H)=o(N_T).
\label{eq:Hprime-H-good}
\end{equation}
The packet endpoint collar defining $V_{\rm end}$ is a coordinate restriction
and does not enter this zero count. Because the perturbation preserves the real
symmetry,
\[
N_{\rm good}(H)=R_{\rm good}(H)+2D(H),
\qquad
N_{\rm good}(H')=R_{\rm good}(H')+2P(H').
\]
Applying Lemma~\ref{lem:real-bookkeeping} on the real interval $[T_1,T_2]$ gives
\[
R_{\rm good}(H')=R_{\rm good}(H)+2B(H)+O(1).
\]
Subtracting the two good-rectangle count identities and using
\eqref{eq:Hprime-H-good} yields
\[
2\{B(H)+P(H')-D(H)\}=o(N_T),
\]
which proves \eqref{eq:D-H}.
\end{proof}

\begin{proof}[Proof of Theorem~\ref{thm:main}]
Fix the good rectangle, and let $D_{\rm good}(F)$ denote the number, with
multiplicity, of reflected nonreal zero pairs of $F$ in its interior. Around
each distinct nonreal zero of $F$ choose a closed disk, disjoint from the real
axis, the other chosen disks, and the boundary. There are finitely many such
disks at fixed $T$. By compactness of their boundaries there is a real
neighbourhood $U_{\rm disk}$ of $(0,0)$ such that
\[
|\eta z+\varepsilon|<|F(z)|
\]
on every disk boundary whenever $(\eta,\varepsilon)\in U_{\rm disk}$.

Choose $H$ and the evaluation-isomorphic packet subfamily by
Lemma~\ref{lem:generic-affine}, with $U_0=U_{\rm disk}$. This fixes the packet
space, $\Gfr_T$, and $\Pret$ before any inertia comparison, and gives
\eqref{eq:perturb-small}.  Hypotheses
\ref{hyp:one-chi-relative-defect} and
\ref{hyp:edge-regular-bounds}, applied to the original Hardy-gauge
matrix $A_F$ on this packet space, give on $\Pret$
\[
A_F=A_T^{\rm true}+E_T^{\rm geom}+R_T^{\rm reg}.
\]
By Proposition~\ref{prop:model-floor},
\[
\lambda_TP_T\succeq a_T\Gfr_T,
\qquad a_T:=0.98\lambda_T.
\]
Apply Lemma~\ref{lem:one-sided-defect-deletion} to
$A_T^{\rm true}$ and $\lambda_TP_T$.  In view of
\eqref{eq:one-chi-relative-defect}, there is a subspace
$V_\beta\subset\Pret$ with
\[
\codim_{\Pret}V_\beta
\le4\beta_-(T;\lambda_T)N_T=O(L^{-2}N_T)
\]
and
\[
A_T^{\rm true}|_{V_\beta}\succeq\frac{a_T}{2}
\Gfr_T|_{V_\beta}.
\]
Put $\Delta_T:=A_H-A_F$.  By \eqref{eq:perturb-small},
\[
(A_T^{\rm true}+\Delta_T)|_{V_\beta}
\succeq\frac{3a_T}{8}\Gfr_T|_{V_\beta}.
\]
Hence, after compression to $V_\beta$,
\[
A_H|_{V_\beta}=(A_T^{\rm true}+\Delta_T)|_{V_\beta}
+E_T^{\rm geom}|_{V_\beta}+R_T^{\rm reg}|_{V_\beta}.
\]
Set
\[
r_T:=\rank(E_T^{\rm geom}|_{V_\beta}),
\qquad
\Gfr_{\beta,T}:=\Gfr_T|_{V_\beta},
\qquad
s_T:=\left\|
\Gfr_{\beta,T}^{-1/2}(R_T^{\rm reg}|_{V_\beta})
\Gfr_{\beta,T}^{-1/2}
\right\|_{S_2}^2.
\]
By \eqref{eq:logonly-edge-rank},
\eqref{eq:logonly-total-relative-HS}, and
\eqref{eq:retained-codim}, one has
$r_T=o(N_T)$, $s_T=o(a_T^2N_T)$, and
$\codim\Pret=o(N_T)$. Apply the ambient-codimension conclusion
\eqref{eq:minmax-codim} of Lemma~\ref{lem:minmax} with the lemma's floor
parameter replaced by
\[
c_T:=\frac{3a_T}{8}>0.
\]
The relative remainder contribution is
$16c_T^{-2}s_T=o(N_T)$.  Including the ambient codimensions of $\Pret$ and
$V_\beta$ therefore gives
\begin{equation}
\boxed{n_-(A_H)=O(L^{-2}N_T)+o(N_T).}
\label{eq:perturbed-inertia}
\end{equation}

The zero-coordinate pullback is exact for $H$, so
\eqref{eq:zero-inertia-generic} and \eqref{eq:perturbed-inertia} give
\[
B(H)+P(H')=O(L^{-2}N_T)+o(N_T).
\]
Lemma~\ref{lem:generic-zero-count} then gives
$D(H)=O(L^{-2}N_T)+o(N_T)$. By the defining
choice $(\eta,\varepsilon)\in U_{\rm disk}$, Rouch\'e's theorem preserves, with
multiplicity, the number of zeros in every chosen disk; all these disks are
disjoint from the real axis. Consequently
\[
D_{\rm good}(F)\le D(H).
\]
Let $D_T(F)$ denote the number, with multiplicity, of reflected nonreal zero
pairs of $F$ with height in $[T,2T]$. Returning from the good rectangle to this
dyadic interval changes the zero count only in two strips of length at most
one. By Lemma~\ref{lem:local-z1-count},
\[
D_T(F)\le D(H)+O(\log T)=O(L^{-2}N_T)+o(N_T).
\]
Because $c>1$, the symmetric rectangle
\eqref{eq:symmetric-contour-rectangle} contains the whole critical strip; at
these heights $\zeta$ has no nontrivial zeros on the complementary half-planes.
Thus every noncritical zero in $[T,2T]$, apart from the endpoint adjustment
already counted above, belongs to exactly one reflected pair. Hence
\[
N_T=N_{0,T}+2D_T(F)+o(N_T),
\]
so
\[
\frac{N_{0,T}}{N_T}=1-O(L^{-2})-o(1)=1-o(1).
\]
We finish by spelling out the dyadic summation. Put
\[
M(T):=N(T)-N_0(T),
\]
and, for a half-open height interval $I$, write
$M(I):=N(I)-N_0(I)$ for the corresponding off-line count.
Use half-open height intervals; by Lemma~\ref{lem:local-z1-count}, changing the
endpoint convention costs $O(\log X)=o(N(X,2X))$. Given $\varepsilon>0$,
choose $X_\varepsilon$ so that for every $X\ge X_\varepsilon$,
\[
M([X,2X))\le\varepsilon N([X,2X)).
\]
Partition $[X_\varepsilon,T)$ into the half-open dyadic blocks
\[
I_j:=\left[\frac{T}{2^{j+1}},\frac{T}{2^j}\right)
\]
whose lower endpoints are at least $X_\varepsilon$; the one remaining bounded
initial segment contributes $O_\varepsilon(1)$. The blocks are disjoint, hence
\[
M(T)
\le\varepsilon\sum_jN(I_j)+O_\varepsilon(1)
\le\varepsilon N(T)+O_\varepsilon(1).
\]
Since $N(T)\to\infty$, division by $N(T)$ gives
\[
\limsup_{T\to\infty}\frac{N(T)-N_0(T)}{N(T)}\le\varepsilon.
\]
Letting $\varepsilon\to0$ proves $N_0(T)/N(T)\to1$.
\end{proof}
